\uuid{8rcF}
\exo7id{3648}
\titre{exo7 3648}
\auteur{quercia}
\organisation{exo7}
\datecreate{2010-03-10}
\isIndication{false}
\isCorrection{false}
\chapitre{Endomorphisme particulier}
\sousChapitre{Autre}
\cpge{mpsi-mp2i.espaces-vectoriels.c, pcsi.espaces-vectoriels.c, ptsi.espaces-vectoriels.c, mpsi-mp2i.dimension-finie.b, pcsi.dimension-finie.b, ptsi.dimension-finie.b, mpsi-mp2i.matrices-applications-lineaires.a, pcsi.matrices-applications-lineaires.a, ptsi.matrices-applications-lineaires.a, bcpst1.espace-vectoriel-kn.b, bcpst2.espaces-vectoriels.b, bcpst2.applications-lineaires-matrices.b}
\module{Algèbre}
\niveau{L2}
\difficulte{}

\contenu{
\question{
Soit $E$ un $ K$-ev de dimension finie $n$,
${\vec u}_1,\dots,{\vec u}_n \in E$ et $f_1,\dots,f_n \in E^*$.

Soit $M$ la matrice de terme général $f_i({\vec u}_j)$.
Montrer que si $\det M \ne 0$, alors $({\vec u}_1,\dots,{\vec u}_n)$
est une base de $E$ et $(f_1,\dots,f_n)$ est une base de $E^*$.
}
}
